\documentclass[11pt,a4paper]{article}
\usepackage[margin=1in]{geometry}
\usepackage{amsmath,amssymb,amsthm}
\usepackage{algorithm}
\usepackage{algpseudocode}
\usepackage{booktabs}
\usepackage{hyperref}

\theoremstyle{definition}
\newtheorem{definition}{Definition}[section]
\newtheorem{theorem}{Theorem}[section]
\newtheorem{lemma}[theorem]{Lemma}
\newtheorem{remark}{Remark}[section]

\title{\textbf{Rigorous Mathematical Specification, Bilinear Multiplicative Interactions, and Convergence of Gated Linear Units (SwiGLU, GeGLU, ReGLU, GLU)}}
\author{Team Scaibu \\ \small Email: \texttt{jaydeep.9664920749@gmail.com} \\ \small Architecture and Core Engine Team}
\date{October 2026}

\begin{document}
\maketitle

\begin{abstract}
This document presents the complete mathematical formulation and algorithmic specification of Gated Linear Units (GLUs), focusing on SwiGLU, GeGLU, ReGLU, and classical GLU. In modern large language models (LLaMA, PaLM, Mistral), standard two-layer feedforward sub-networks are replaced by three-matrix bilinear gated projections: $\operatorname{FFN}_{\text{SwiGLU}}(x) = (\operatorname{SiLU}(x W_{\text{gate}}) \odot x W_{\text{up}}) W_{\text{down}}$. We formalize bilinear gating theory, provide the complete algorithmic pseudo-code matching the reference implementation, derive exact computational complexities, step through a numerical example, and analyze dimension scaling.
\end{abstract}

\section{Notation and Mathematical Preliminaries}

Let $X \in \mathbb{R}^{B \times d_{\text{in}}}$ denote an input mini-batch tensor across $B$ tokens and $d_{\text{in}}$ embedding channels.

\begin{itemize}
    \item $W_{\text{gate}} \in \mathbb{R}^{d_{\text{hid}} \times d_{\text{in}}}$: Gate projection weight matrix.
    \item $W_{\text{up}} \in \mathbb{R}^{d_{\text{hid}} \times d_{\text{in}}}$: Value (up) projection weight matrix.
    \item $W_{\text{down}} \in \mathbb{R}^{d_{\text{out}} \times d_{\text{hid}}}$: Optional output down-projection weight matrix.
    \item $\sigma$: Non-linear gating activation (e.g., SiLU for SwiGLU, GELU for GeGLU, Sigmoid for GLU).
    \item $\odot$: Element-wise Hadamard product operator.
\end{itemize}

\begin{table}[h]
\centering
\caption{Summary of Mathematical Symbols for Gated Linear Units}
\begin{tabular}{lll}
\toprule
\textbf{Symbol} & \textbf{Type / Domain} & \textbf{Description} \\
\midrule
$X$ & $\mathbb{R}^{B \times d_{\text{in}}}$ & Input embedding batch \\
$W_{\text{gate}}, W_{\text{up}}$ & $\mathbb{R}^{d_{\text{hid}} \times d_{\text{in}}}$ & Gating and value projection matrices \\
$W_{\text{down}}$ & $\mathbb{R}^{d_{\text{out}} \times d_{\text{hid}}}$ & Down projection matrix \\
$\odot$ & Operator & Hadamard element-wise product \\
\bottomrule
\end{tabular}
\end{table}

\section{Problem Definition and Core Concepts}

\subsection{Intuition and Motivation}
Standard feedforward layers compute $y = \sigma(x W_1) W_2$. In GLU variants, the representation is split into two parallel paths: a linear value channel $x W_{\text{up}}$ and a non-linear gate channel $\sigma(x W_{\text{gate}})$. The multiplicative interaction acts as a continuous coordinate filter: when the gate output is close to 0, information along that channel is suppressed; when 1, it passes through. Shazeer (2020) demonstrated that SwiGLU consistently outperforms ReLU/GELU feedforward blocks in transformer language models.

\subsection{Formal Mathematical Definition}
\begin{definition}[GLU Family Formulations]
For input vector $x \in \mathbb{R}^{d_{\text{in}}}$:
\begin{align}
\operatorname{GLU}(x) &= \sigma\left(x W_{\text{gate}}^T\right) \odot \left(x W_{\text{up}}^T\right) \label{eq:glu_classic} \\
\operatorname{SwiGLU}(x) &= \operatorname{SiLU}\left(x W_{\text{gate}}^T\right) \odot \left(x W_{\text{up}}^T\right) \label{eq:swiglu} \\
\operatorname{GeGLU}(x) &= \operatorname{GELU}\left(x W_{\text{gate}}^T\right) \odot \left(x W_{\text{up}}^T\right) \label{eq:geglu} \\
\operatorname{ReGLU}(x) &= \operatorname{ReLU}\left(x W_{\text{gate}}^T\right) \odot \left(x W_{\text{up}}^T\right) \label{eq:reglu}
\end{align}
When combined with down-projection:
\begin{equation}
\operatorname{FFN}_{\text{GLU}}(x) = \left( \sigma\left(x W_{\text{gate}}^T\right) \odot \left(x W_{\text{up}}^T\right) \right) W_{\text{down}}^T \label{eq:ffn_glu}
\end{equation}
\end{definition}

\section{Algorithmic Specification}

The computational pipeline implemented in \texttt{NnAlgoGatedLinearUnits} is detailed below.

\begin{algorithm}[H]
\caption{Gated Linear Units (SwiGLU/GeGLU/ReGLU) Forward Evaluation}
\begin{algorithmic}[1]
\Require Input batch $X \in \mathbb{R}^{B \times d_{\text{in}}}$, weights $W_{\text{gate}}, W_{\text{up}} \in \mathbb{R}^{d_{\text{hid}} \times d_{\text{in}}}$, optional $W_{\text{down}} \in \mathbb{R}^{d_{\text{out}} \times d_{\text{hid}}}$
\Require Variant $\in \{\text{swiglu}, \text{geglu}, \text{reglu}, \text{glu}\}$
\Ensure Output batch $Y$, hidden activations $H \in \mathbb{R}^{B \times d_{\text{hid}}}$
\State $B \gets \text{rows}(X), \quad d_{\text{in}} \gets \text{cols}(X), \quad d_{\text{hid}} \gets \text{rows}(W_{\text{gate}})$
\State $G \gets X \cdot W_{\text{gate}}^T \in \mathbb{R}^{B \times d_{\text{hid}}}$
\State $U \gets X \cdot W_{\text{up}}^T \in \mathbb{R}^{B \times d_{\text{hid}}}$
\State $H \gets \text{new Matrix of shape } (B, d_{\text{hid}})$
\For{$b \gets 1 \textbf{ to } B$}
    \For{$j \gets 1 \textbf{ to } d_{\text{hid}}$}
        \State $g \gets G[b][j], \quad u \gets U[b][j]$
        \If{$\text{variant} = \text{swiglu}$}
            \State $\text{act} \gets g \cdot (1.0 / (1.0 + \exp(-g)))$
        \ElsIf{$\text{variant} = \text{geglu}$}
            \State $\text{act} \gets 0.5 \cdot g \cdot (1.0 + \operatorname{erf}(g / \sqrt{2.0}))$
        \ElsIf{$\text{variant} = \text{reglu}$}
            \State $\text{act} \gets \max(0.0, g)$
        \ElsIf{$\text{variant} = \text{glu}$}
            \State $\text{act} \gets 1.0 / (1.0 + \exp(-g))$
        \EndIf
        \State $H[b][j] \gets \text{act} \cdot u$
    \EndFor
\EndFor
\If{$W_{\text{down}} \text{ is not null}$}
    \State $Y \gets H \cdot W_{\text{down}}^T \in \mathbb{R}^{B \times d_{\text{out}}}$
\Else
    \State $Y \gets H$
\EndIf
\State \Return $\{\text{output\_batch}: Y, \text{hidden\_gated}: H\}$
\end{algorithmic}
\end{algorithm}

\section{Theoretical Guarantees and Parameter Parity}

\begin{theorem}[Parameter and Compute Parity Scaling Rule (Shazeer 2020)]
A standard 2-layer FFN with hidden dimension $d_{\text{ffn}} = 4 d_{\text{in}}$ contains $8 d_{\text{in}}^2$ parameters. To maintain exact parameter and FLOP parity in a 3-matrix SwiGLU FFN ($3 d_{\text{in}} d_{\text{hid}} = 8 d_{\text{in}}^2$), the hidden dimension must be scaled to:
\begin{equation}
d_{\text{hid}} = \left\lfloor \frac{8}{3} d_{\text{in}} \right\rfloor \approx \frac{2}{3} \cdot 4 d_{\text{in}}
\end{equation}
\end{theorem}

\subsection{Limitations}
\begin{itemize}
    \item \textbf{Three Matrix Multiplications}: Requires 3 separate projection operations per token, which can be memory bandwidth bound on GPUs without tensor fusion kernels.
\end{itemize}

\section{Complexity Analysis}

\subsection{Time Complexity}
\begin{itemize}
    \item \textbf{Gate and Up Projections}: $2 \times (B \cdot d_{\text{in}} \cdot d_{\text{hid}})$ FLOPs.
    \item \textbf{Hadamard Product}: $\mathcal{O}(B \cdot d_{\text{hid}})$ FLOPs.
    \item \textbf{Down Projection}: $B \cdot d_{\text{hid}} \cdot d_{\text{out}}$ FLOPs.
    \item \textbf{Total Time}: $\Theta\left(B \cdot d_{\text{hid}} (2 d_{\text{in}} + d_{\text{out}})\right)$ FLOPs.
\end{itemize}

\subsection{Space Complexity}
\begin{itemize}
    \item \textbf{Auxiliary Memory}: Stores gate buffer $G$, up buffer $U$, and gated product $H$ of shape $B \times d_{\text{hid}}$ ($\Theta(B \cdot d_{\text{hid}})$ space).
\end{itemize}

\section{Worked Numerical Example}

Consider $B=1, d_{\text{in}}=2, d_{\text{hid}}=2$, $\text{variant} = \text{swiglu}$:
$x = [1.0, 0.0]$, $W_{\text{gate}} = \begin{bmatrix} 1.0 & 0.0 \\ -1.0 & 0.0 \end{bmatrix}$, $W_{\text{up}} = \begin{bmatrix} 2.0 & 0.0 \\ 3.0 & 0.0 \end{bmatrix}$.

\begin{enumerate}
    \item $G = x W_{\text{gate}}^T = [1.0(1) + 0, 1.0(-1) + 0] = [1.0, -1.0]$.
    \item $U = x W_{\text{up}}^T = [1.0(2) + 0, 1.0(3) + 0] = [2.0, 3.0]$.
    \item \textbf{SiLU on Gate}:
    \begin{align}
    \text{SiLU}(1.0) &= 1.0 \times \frac{1}{1 + e^{-1}} \approx 0.731059 \\
    \text{SiLU}(-1.0) &= -1.0 \times \frac{1}{1 + e^1} \approx -0.268941
    \end{align}
    \item \textbf{Hadamard Product}:
    \begin{align}
    H_1 &= 0.731059 \times 2.0 \approx 1.462118 \\
    H_2 &= -0.268941 \times 3.0 \approx -0.806823
    \end{align}
\end{enumerate}

\section{Numerical Considerations and Edge Cases}

\begin{itemize}
    \item \textbf{Fused Gate and Up GEMM}: $W_{\text{gate}}$ and $W_{\text{up}}$ are typically stacked horizontally into a single $(2 d_{\text{hid}} \times d_{\text{in}})$ matrix, executing one combined GEMM.
\end{itemize}

\begin{thebibliography}{9}
\bibitem{shazeer2020} N.~Shazeer, ``GLU variants improve transformer,'' \emph{arXiv preprint arXiv:2002.05202}, 2020.
\bibitem{dauphin2017} Y.~N.~Dauphin et al., ``Language modeling with gated convolutional networks,'' in \emph{ICML}, pp.~933--941, 2017.
\bibitem{touvron2023llama1} H.~Touvron et al., ``Llama: Open and efficient foundation language models,'' \emph{arXiv:2302.13971}, 2023.
\end{thebibliography}

\end{document}
